\section{The actual ideal image of a quasi-coherent module morphism}
Let $X$ be any scheme. Write $\mathcal O_X$ for its original structure module.
Let $M$ be an actual quasi-coherent $\mathcal O_X$-module and let
$\varphi:M\to\mathcal O_X$ be an actual module-sheaf morphism.
For an affine open $U$, define
\[
 I_\varphi(U)=\operatorname{im}(\Gamma(U,\varphi))\subseteq\Gamma(U,\mathcal O_X).
\]
The module structure on the codomain is the original ring acting on itself;
thus this linear image is literally an ideal of the original section ring.
There are no field, finiteness, reducedness, separation, or nonemptiness hypotheses.

\begin{lemma}[The original affine section image localizes]
\label{lem:quasicoherent-image-ideal-localization}
\leanok
\lean{EnvelopeProperty.QuasicoherentImageIdeal.image_localization}
\uses{prop:actual-module-image-ideal-sheaf}
For every affine $U$ and original $f\in\Gamma(U,\mathcal O_X)$, with
$r:\Gamma(U,\mathcal O_X)\to\Gamma(D_U(f),\mathcal O_X)$ the actual restriction,
\[
 r(I_\varphi(U))\,\Gamma(D_U(f),\mathcal O_X)=I_\varphi(D_U(f)).
\]
\end{lemma}
\begin{proof}
Apply the existing whole basic-open image-localization theorem
attached to Proposition~\ref{prop:actual-module-image-ideal-sheaf},
with the SAME original $X$, $M$, $\varphi$, affine $U$ and function $f$.
Its affine ideal is defined to be the range of the original section map.
Unfolding that definition gives exactly the displayed equality, with
no transport of the original ring, module, morphism or affine open.
\end{proof}

\begin{theorem}[The genuine ideal sheaf image]
\label{thm:quasicoherent-image-ideal}
\leanok
\lean{EnvelopeProperty.QuasicoherentImageIdeal.imageIdeal,EnvelopeProperty.QuasicoherentImageIdeal.mem_imageIdeal}
\uses{lem:quasicoherent-image-ideal-localization}
The actual ideals $I_\varphi(U)$ form an original $X.\mathrm{IdealSheafData}$,
whose affine ideal at every original affine $U$ is exactly
$\operatorname{im}(\Gamma(U,\varphi))$. In particular, membership means
existence of an actual section $m\in\Gamma(U,M)$ mapping to the given
original regular function.
\end{theorem}
\begin{proof}
Use these original affine image ideals as the ideal fields in the existing
ideal-sheaf-data constructor. Its required basic-open equality is precisely
Lemma~\ref{lem:quasicoherent-image-ideal-localization}. Use the constructor's
default intersection of actual affine zero loci for its support; there is no
additional support identification hypothesis. The affine formula is the
construction itself, and membership is the definition of the linear-map range.
\end{proof}

This is a lower geometric producer, not a relative-canonical discrepancy theorem.
Applied subsequently to the genuine evaluation map from the dual of an existing
line sheaf, it will give the actual zero ideal of an original section.
Deriving Cartier regularity from nonzero sections on an integral scheme,
deriving nonvanishing of the original canonical differential for birational maps,
and comparing its divisor with intrinsic DVR discrepancy remain separate obligations.
Uniform multiplier global generation is not claimed here.
